\documentclass[11pt,letterpaper]{article}
\usepackage[margin=1in]{geometry}
\usepackage{titlesec}
\usepackage{enumitem}
\usepackage{xcolor}
\usepackage{booktabs}
\usepackage{hyperref}
\usepackage{fancyhdr}
\hypersetup{colorlinks=true,urlcolor=blue!60!black,linkcolor=black}
\titleformat{\section}{\Large\bfseries}{\thesection}{0.6em}{}
\titleformat{\subsection}{\large\bfseries}{\thesubsection}{0.6em}{}
\setlength{\parskip}{4pt}\setlength{\parindent}{0pt}
\pagestyle{fancy}\fancyhf{}
\rhead{\small Hanzo console / Langfuse fork --- tRPC migration scoping}
\rfoot{\small \thepage}
\title{\textbf{console (Langfuse fork): tRPC $\rightarrow$ @zap-proto/web migration decision}}
\author{Architectural scoping --- read-only survey}
\date{}

\begin{document}
\maketitle

\section*{TL;DR --- Recommendation}
\textbf{Option B: keep console on tRPC; expose a single tRPC$\leftrightarrow$ZAP gateway at the cluster edge.} The existing \texttt{ui-customization} migration already proves the gateway pattern works in-process. Migrating all 74 routers permanently inherits Langfuse's churn ($\approx$805 router-touching upstream commits in the last 9 months) as merge conflicts on every sync; the engineering cost dwarfs the architectural-purity benefit, and the ``native .zap only'' directive is satisfied for the rest of the stack with one bridge endpoint instead of 74.

\section{Scope of the surface}
\begin{itemize}[leftmargin=1.2em,itemsep=2pt]
  \item \textbf{Routers:} 74 files containing \texttt{createTRPCRouter} under \texttt{web/src}.
  \item \textbf{Procedures:} $\approx$342 procedures across those routers.
  \item \textbf{Frontend call-sites:} $\approx$2{,}133 \texttt{api.<router>.<proc>.use*} call-sites across \texttt{web/src}. Every one is a code change.
  \item \textbf{tRPC deps:} \texttt{@trpc/client}, \texttt{@trpc/next}, \texttt{@trpc/react-query}, \texttt{@trpc/server} pinned to \texttt{\^{}11.13.4}.
  \item \textbf{Largest single router:} \texttt{features/datasets/server/dataset-router.ts} --- 1{,}933 LOC, 39 procedures, Prisma + ClickHouse + S3 + batch jobs.
  \item \textbf{Fork posture:} \texttt{git remote -v} shows \texttt{upstream $\rightarrow$ langfuse/langfuse.git}; merge commits like \texttt{34b0323a3 merge: upstream/main} confirm an active sync cadence.
  \item \textbf{Branch:} already on \texttt{zap-native-ts-stack}.
\end{itemize}

\section{Upstream churn (the load-bearing number)}
Measured against \texttt{upstream/main}:
\begin{center}
\begin{tabular}{lr}
\toprule
Window & Router-touching commits \\
\midrule
9 months ending Q2 2026 & 805 \\
5.5 months ending Q2 2026 & 503 \\
Q2 2026 alone & 278 \\
\bottomrule
\end{tabular}
\end{center}
Hot routers in the last 12 months: \texttt{dataset-router.ts} 57 touches, \texttt{traces.ts} 37, \texttt{evals/router.ts} 31, \texttt{scores.ts} 29, \texttt{sessions.ts} 22, \texttt{promptRouter.ts} 22. These are the same routers that drive primary user flows --- migrating them away from tRPC creates the worst-case ongoing merge cost.

\section{The proof point: \texttt{ui-customization}}
The only router migrated so far (\texttt{a6c285ee5 feat(zap): native ZAP capability-RPC foundation + migrate ui-customization router}) is the easiest possible case: $\approx$70 LoC, two read-only methods, no DB writes. Its migration shipped:
\begin{itemize}[leftmargin=1.2em,itemsep=2pt]
  \item A new Go service binary as the backend (no in-process server).
  \item A hand-written \texttt{web/proto/ui-customization.zap} schema.
  \item Hand-checked-in \texttt{web/src/server/zap/gen/ui-customization\_zap.ts} (no \texttt{zapgen} in the TS build).
  \item A bridge at \texttt{web/src/pages/api/zap/[[...path]].ts} (HTTP $\rightarrow$ TCP $\rightarrow$ \texttt{@hanzo/zap}).
  \item A custom \texttt{web/src/utils/zapClient.ts} that replaces \texttt{api.uiCustomization.get.useQuery()}.
  \item A dev-only capability minter (the canonical ML-DSA-65 \texttt{Capability} is not yet shipped --- \texttt{@zap-proto/web v0.1.1} marks \texttt{mintCap} as a slot for ``a later version'').
\end{itemize}
Extrapolating this template to \texttt{dataset-router.ts} (39 procedures across Prisma, ClickHouse, S3, Stripe, BullMQ) is not a days-of-work exercise --- it is a multi-week effort \emph{per hot router}.

\section{Option A --- Full migration to @zap-proto/web}
\textbf{Cost.} Three router classes; midpoint estimates:
\begin{center}
\begin{tabular}{lrrr}
\toprule
Class & Count & Wks each & Total \\
\midrule
Trivial CRUD & 25 & 0.5 & 12.5 \\
Medium & 35 & 1.0 & 35.0 \\
Heavy & 14 & 2.0 & 28.0 \\
\midrule
\textbf{Subtotal} & 74 & --- & \textbf{75.5} \\
Cross-cutting (zapgen wiring, capability mint, auth, tests, rollout) & --- & --- & 10--20 \\
\midrule
\textbf{Total Option A} & & & \textbf{$\approx$85--95 eng-weeks} \\
\bottomrule
\end{tabular}
\end{center}

\textbf{Risk hotspots.}
\begin{itemize}[leftmargin=1.2em,itemsep=2pt]
  \item Permanent merge tax: every one of the $\approx$800 router-commits/yr upstream becomes a manual rebase against a fork. Annualized: $\approx$8--12 eng-weeks/yr of ongoing merge labor.
  \item \texttt{@zap-proto/web} is \texttt{v0.1.1}, pre-1.0; the production-grade \texttt{Capability} mint (ML-DSA-65, BLAKE3) is not shipped. Migrating EE billing/SSO routers before that lands is unsafe.
  \item EE features (\texttt{verified-domains}, \texttt{multi-tenant-sso}, \texttt{cloudBilling}, \texttt{spend-alerts}, \texttt{sso-config}) span 5 routers with audit-sensitive logic and tight coupling to NextAuth and Stripe.
  \item No \texttt{zapgen --target=ts} step in console's turbo graph today; CI/regen pipeline is a prerequisite.
  \item Frontend churn: $\approx$2{,}133 call-sites to rewrite, plus React Query cache-key compatibility decisions.
\end{itemize}

\textbf{Operational impact.} Can be done router-by-router (the gateway approach makes mixed deployment easy), but full cutover requires removing \texttt{@trpc/*} from \texttt{package.json}, which is a one-shot all-or-nothing landing.

\section{Option B --- Keep tRPC; bridge at the edge}
\textbf{Shape.} A single Go (or Node) service in front of the cluster terminates ZAP on one side and speaks tRPC HTTP-batch to console on the other. tRPC's wire format is already JSON over HTTP --- the bridge is a few hundred lines of glue: it accepts a ZAP envelope with method ordinal + capability, validates the cap, maps ordinal $\rightarrow$ tRPC procedure path, replays the JSON body to the console URL, and re-envelopes the JSON reply.

\textbf{Cost.}
\begin{itemize}[leftmargin=1.2em,itemsep=2pt]
  \item Gateway implementation: 2--3 eng-weeks (Go service or Node sidecar; reuses \texttt{@hanzo/zap} TCP server).
  \item Capability $\rightarrow$ NextAuth-session shim: 1 eng-week.
  \item One \texttt{console.zap} schema enumerating the externally-callable procedures: 1 eng-week.
  \item \textbf{Total upfront: $\approx$4--5 eng-weeks.}
  \item Ongoing: $\approx$0.5--1 eng-week/quarter to keep the schema in sync.
\end{itemize}

\textbf{Risk hotspots.}
\begin{itemize}[leftmargin=1.2em,itemsep=2pt]
  \item ``A bridge counts as sludge.'' The directive says ``no proxy no hacks no sludge native .zap only.'' Mitigation: the bridge is \emph{one} component, owned, with a tested schema --- not a sprinkling of compat shims.
  \item Capability semantics don't map 1:1 to NextAuth session scopes; the shim enforces permission caveats at the gateway, not in console.
  \item If upstream Langfuse one day moves off tRPC, the bridge changes shape too. Probability is low in the near term.
\end{itemize}

\textbf{Operational impact.} Zero console downtime. The bridge is deployed alongside the existing console pod; clients are flipped one feature at a time via the dispatch table.

\section{Option C --- Greenfield Hanzo console on ZAP}
\textbf{Cost.} Feature parity for ``what Hanzo actually uses'' ($\approx$12 of the 74 routers):
\begin{itemize}[leftmargin=1.2em,itemsep=2pt]
  \item Backend services (Go + ZAP): 12--16 eng-weeks.
  \item Frontend (Next.js, charts, tables, batch tooling): 25--40 eng-weeks.
  \item Auth, RBAC, audit logging, telemetry: 5--8 eng-weeks.
  \item \textbf{Total Option C (subset): $\approx$45--60 eng-weeks.}
  \item Full Langfuse parity: \textbf{$\approx$120+ eng-weeks}.
\end{itemize}

\textbf{Risk hotspots.} Loses every feature Langfuse upstream ships next, forever. Forfeits a $>$100-eng-yr open-source asset.

\section{Recommendation and justification}
\textbf{Adopt Option B.} The gateway is $\approx$15--20$\times$ cheaper than Option A upfront, eliminates the permanent fork-merge tax (the dominant lifetime cost), and preserves Langfuse's free upstream velocity. The architectural directive --- ``native .zap only'' --- is satisfied at the boundary that actually matters: every other Hanzo service speaks ZAP to console through the bridge; console is the only consumer of console's tRPC. The directive is not violated by having tRPC live \emph{inside} one fork; it would be violated by having tRPC leak \emph{between} services. The bridge enforces exactly that boundary.

If, in 12--18 months, \texttt{@zap-proto/web} hits 1.0 with shipped \texttt{Capability} mint and \texttt{zapgen --target=ts} is wired into console's build, revisit Option A starting with one cold router (e.g. \texttt{auditLogs} or \texttt{notificationPreferences}) as a fully-native proof-of-concept --- but do not commit to migrating hot routers until upstream Langfuse stabilizes or Hanzo decides to detach from the upstream entirely.

\section{If migration is chosen later --- phased plan}
\begin{description}[leftmargin=1.4em,itemsep=2pt]
  \item[Phase 0 (foundation, 3--4 wk).] Wire \texttt{zapgen --target=ts} into turbo; add a \texttt{web/proto/} convention and CI regen check; ship the production capability mint (depends on \texttt{@zap-proto/web} 1.0); add a \texttt{zapProcedure} adapter that gives migrated routers a tRPC-shaped DX from the codegen'd views.
  \item[Phase 1 (cold-path routers, 6--8 wk).] Migrate \texttt{auditLogs}, \texttt{notificationPreferences}, \texttt{tableViewPresets}, \texttt{dashboardWidgets}, \texttt{surveys}, \texttt{userAccount}, \texttt{utilities}. No upstream churn risk; all low-cardinality reads.
  \item[Phase 2 (hot reads, 10--14 wk).] \texttt{traces}, \texttt{observations}, \texttt{sessions}, \texttt{scores}, \texttt{public}. Coordinate every merge with a frozen \texttt{.zap} schema.
  \item[Phase 3 (writes \& EE, 16--20 wk).] \texttt{datasets}, \texttt{prompts}, \texttt{evals}, \texttt{experiments}, plus EE: \texttt{cloudBilling}, \texttt{ssoConfig}, \texttt{verifiedDomain}, \texttt{spendAlert}.
  \item[Phase 4 (cutover, 3--4 wk).] Remove \texttt{@trpc/*} from \texttt{package.json}; delete \texttt{web/src/server/api/trpc.ts} and \texttt{root.ts}; declare upstream divergence formal.
\end{description}

\section{Critical files for any chosen path}
\begin{itemize}[leftmargin=1.2em,itemsep=2pt]
  \item \texttt{web/src/server/api/trpc.ts} --- procedure builders, auth context, the single chokepoint.
  \item \texttt{web/src/server/api/root.ts} --- the canonical router registry; the migration ledger.
  \item \texttt{web/src/pages/api/zap/[[...path]].ts} --- existing ZAP bridge; template for both Option A and Option B.
  \item \texttt{web/src/utils/zapClient.ts} --- existing client surface; becomes the gateway client.
  \item \texttt{web/src/features/datasets/server/dataset-router.ts} --- the canary for ``can we afford Option A?''.
\end{itemize}

\end{document}
